% TOP_PROBLEM: {"id": 3079, "title": "Kneser–Poulsen conjecture", "category_id": 6, "category": {"id": 6, "name": "geometry", "display_name": "Geometry", "description": "Euclidean and non-Euclidean geometry, geometric structures.", "slug": "geometry", "order_index": 6, "created_at": "2026-07-31T15:26:25.671Z"}, "status": "open", "problem_number": "OPG-2089", "set_id": 12}
% FIRST_POSTED: 2026-10-01
% ENTRY_KIND: statement_only
% UnsolvedMath ID 3079; source catalogue code OPG-2089
% !TeX program = lualatex
% Definition and sources only; this file is not an LLM solution attempt.
\documentclass[11pt,a4paper]{article}
\usepackage[margin=25mm]{geometry}
\usepackage{fontspec}
\setmainfont{lmroman10-regular.otf}[BoldFont=lmroman10-bold.otf,
 ItalicFont=lmroman10-italic.otf,BoldItalicFont=lmroman10-bolditalic.otf]
\usepackage{amsmath,amssymb,mathtools}
\usepackage{unicode-math}
\setmathfont{latinmodern-math.otf}
\AtBeginDocument{\let\setminus\smallsetminus}
\usepackage{xurl,hyperref}
\hypersetup{colorlinks=true,urlcolor=blue}
\setcounter{secnumdepth}{0}
\setlength{\parindent}{0pt}
\setlength{\parskip}{5pt}
\setlength{\emergencystretch}{3em}
\begin{document}
\section*{Kneser–Poulsen conjecture}\label{3079}
{\small UnsolvedMath ID 3079 · OPG-2089 · Geometry · OpenGarden}
\subsection{Definitions and mathematical statement}
Let $n,N\ge1$ and consider two labelled configurations of centres
$p_1,\ldots,p_N$ and $q_1,\ldots,q_N$ in $\mathbb R^n$.
Repeated centres are allowed. For $x\in\mathbb R^n$, let
$B(x,1)=\{z:\|z-x\|_2\le1\}$ be the closed unit ball, and let
$\operatorname{vol}_n$ denote $n$-dimensional Lebesgue measure.
Assume the second configuration is an expansion of the first in the
pairwise-distance sense:
\[
 \|q_i-q_j\|_2\ge\|p_i-p_j\|_2
                    \qquad(1\le i,j\le N).
\]
The conjecture asserts
\[
 \operatorname{vol}_n\!\left(\bigcup_{i=1}^N B(q_i,1)\right)
 \ge
 \operatorname{vol}_n\!\left(\bigcup_{i=1}^N B(p_i,1)\right).
\]
The volume counts overlapping regions once, as ordinary volume of a
union, rather than as a sum of ball volumes. The same radius is assigned
to every ball both before and after the rearrangement; replacing one by
any other common positive radius gives an equivalent assertion by scaling.

Only the two endpoint configurations must satisfy the displayed distance
inequalities. The problem does not assume that they can be joined by a
continuous motion during which every pairwise distance is nondecreasing.
Such a hypothesis would restrict the class of rearrangements. Both
configurations must lie in the same specified Euclidean dimension.
The source entry asks about the union of equal balls; corresponding
questions for unequal radii or intersections are distinct extensions.
\subsection{Short English statement}
Can expanding all pairwise distances between centres ever decrease the volume covered by equal balls?
\subsection{Sources}
\begin{itemize}
\item[S1] ULAM.AI and the UnsolvedMath Contributors, supplied problems.json, record 3079 (OPG-2089), OpenGarden collection. Catalogue identity and source transcription; CC BY 4.0.
\url{https://huggingface.co/datasets/ulamai/UnsolvedMath}.
\item[S2] Open Problem Garden, Kneser–Poulsen conjecture. Original problem and discussion; the mathematical formulation below is expanded from the supplied source record.
\url{http://www.openproblemgarden.org/op/kneser_poulsen_conjecture}.
\end{itemize}
\end{document}
